% S4: ban TIENG VIET cua paper2/blocks/08_subsidence.tex.
% Chi duoc \input boi paper2_hujos/supplementary_vi.tex.
% NGUYEN TAC: giu nguyen moi con so va moi hedge; thap phan kieu Viet ({,});
% tham chieu cheo toi bai chinh dung \eqref/\ref (so lay tu main.aux
% boi experiments/make_supplementary_stubs.py) thay vi go tay "Eq.~(16)" — ban EN go tay
% nen da lech so voi ban VI (eq:umax EN=16, VI=15).
\section*{S4. Tính bất biến đối với sụt lún}

Tuân thủ AWD thường được phát biểu bằng một mực nước ngầm hoặc độ sâu ngập tuyệt
đối, ví dụ $15$\,cm dưới mặt ruộng, tức ngưỡng tưới lại của nghiệm thức Safe-AWD
trong~\cite{balaine2019}. Ở một đồng bằng đang sụt lún, cao độ mặt ruộng $z_s(t)$ hạ
với tốc độ $\dot z$. Hai giao thức đo cho hai kết quả khác nhau và cần tách bạch:
nếu ngưỡng được quy về một \emph{mốc cao độ cố định} thì mốc và mặt ruộng rời nhau
nên ngưỡng hữu hiệu trôi; còn nếu ngưỡng được đo \emph{tương đối với mặt ruộng hiện
tại}, như ống đo mực nước perched đặt theo mặt ruộng, thì cả mốc lẫn mặt ruộng cùng
hạ nên ngưỡng không trôi. Mệnh đề dưới đây chỉ áp dụng cho giao thức thứ nhất.

\begin{proposition}[Độ trôi của một ngưỡng neo theo mốc cố định]
\label{prop:subsid}
Xét một ngưỡng được biểu diễn theo một \emph{cao độ gốc cố định}, chẳng hạn mực
nước đọc từ một cảm biến neo theo mốc trắc đạc. Trong một chu kỳ cấp tín chỉ dài
$T_c$ năm, ngưỡng đó dịch đi $\dot z\,T_c$ so với mặt ruộng. Với
$\dot z\in[1{,}1;2{,}5]$~cm/năm và $T_c=5$~năm, độ dịch là $5{,}5$--$12{,}5$~cm,
tức $37$--$83\%$ của một ngưỡng $15$\,cm. Một tiêu chí viết theo tỷ lệ cạn kiệt
\eqref{eq:depl} có độ trôi bằng không trong cả hai giao thức, vì $FC$ và $WP$ là
tính chất của đất quy về tầng rễ chứ không quy về một cao độ tuyệt đối.
\end{proposition}

Phạm vi của mệnh đề cần được nói rõ. Kết luận trên \emph{không} áp dụng cho một
ngưỡng đo tương đối với mặt ruộng hiện tại: tiêu chí $15$\,cm dưới mặt ruộng của
nghiệm thức Safe-AWD trong~\cite{balaine2019} là một mực nước perched đo từ mặt
ruộng, nên nó tịnh tiến cùng mặt ruộng và không tự động trôi hết lượng sụt lún.
Chúng tôi không xác minh được từ nguồn đó loại ống đo mà các điểm triển khai dùng,
nên không gán giao thức nào cho chúng; độ dịch $5{,}5$--$12{,}5$~cm chỉ là cận trên
đúng cho trường hợp cảm biến neo theo mốc cố định.

\begin{proof}
Tiêu chí neo theo mốc so một độ sâu đo được với một mốc không đổi; nếu mốc và mặt
ruộng rời nhau với tốc độ $\dot z$ thì ngưỡng hữu hiệu thay đổi $\dot z T_c$ sau
$T_c$ năm. Nếu mốc được đặt lại theo mặt ruộng hiện tại thì hai đại lượng tịnh tiến
cùng nhau và độ trôi bằng không; đó là lý do mệnh đề chỉ phát biểu cho giao thức
thứ nhất.
Tiêu chí cạn kiệt so $FC-w_t$ với $(FC-WP)$, cả hai đều được định nghĩa bên trong tầng
rễ và tịnh tiến cứng cùng mặt ruộng, nên tỷ số của chúng bất biến đối với $z_s$.
\end{proof}

Các tốc độ sụt lún là số đo cho đồng bằng này trong~\cite{minderhoud2017}. Hệ quả thực
tiễn là một phương án neo theo tỷ lệ cạn kiệt không cần đo đạc lại các mốc ruộng trong
suốt chu kỳ cấp tín chỉ năm năm, còn phương án dùng độ sâu tuyệt đối thì cần, với điều kiện phân biệt ở trên:
nếu đọc từ cảm biến neo theo mốc cố định thì mốc cao độ đó phải được đo đạc lại
theo chu kỳ, còn nếu đo tương đối mặt ruộng thì không. Trong cả hai trường hợp,
tỷ lệ cạn kiệt vẫn là đại lượng không cần hiệu chỉnh theo cao độ.
